% ======================================================================
%  Appendix A: the model, the predictor, and the proofs
%  Sections A1-A3 are new; A4-A7 carry the protocols, gates, lemmas and
%  the status table that were already in this appendix.
% ======================================================================
\section{The model, the predictor, and the proofs}\label{app:evidence}

This appendix does three things the main text does not have room for. It fixes the model
formally (A1), including which of the two readings of $\varphi$ is in force at each point --- a
distinction the text uses throughout without naming. It defines the predictor and the quantity
$\kappa$ that the classification is stated in (A2). It then proves, or cites the verification of,
the structural facts the paper relies on (A4), and records the protocols (A5), the gates (A6) and
the status of each claim (A7).

\subsection{The model}\label{app:model}

\begin{definition}[Rule]\label{def:rule}
A \emph{rule} is a pair $\rho=(\varphi,f)$ with
\[
\varphi:W\to A,\qquad W=\bigcup_{m\ge1}\Gamma^{m},\qquad A\ \text{finite},
\]
\[
f:A\times\mathbb N\times\mathbb N^{4}\longrightarrow
\{\ell{+}1,\ \ell{-}1,\ r{+}1,\ r{-}1,\ \text{left trim},\ \text{right trim},\ \text{double trim},\ \text{halt}\},
\]
evaluated at $\bigl(\varphi(\tau[\ell..r]),\,n,\,\ell,\,r,\,L,\,R\bigr)$. Two conditions hold by
definition and are used silently below: $\varphi$ depends on the current window only, never on the
history, so $(\varphi,f)$ makes the state $(L,R,\ell,r)$ Markov; and every convention by which a
rule resolves an inadmissible move, a trim or a halt is part of $f$. No computability, complexity or
measurability condition is placed on $\varphi$: it is the \emph{specification} of the rule, not an
object the rule computes.
\end{definition}

\begin{definition}[Run]\label{def:run}
Fix $n$ and $\tau\in\Gamma^{n}$. The state is $(L,R,\ell,r)$ with $0\le L\le\ell\le r\le R\le n-1$,
initially $(0,n-1,0,n-1)$. Given a state, let $a=f(\varphi(\tau[\ell..r]),n,\ell,r,L,R)$.
\begin{enumerate}
\item If $\ell=r$, the run halts and outputs $\delta=\ell/(n-1)$.
\item If $a$ is a move taking a header to $p$, it is admissible iff $L\le p\le R$. If admissible and
$p$ meets or passes the other header, the run halts and outputs $\delta=p/(n-1)$; if inadmissible the
run aborts with no output.
\item If $a$ is a trim, it is admissible iff the new region is non-empty; then $(L,R)\leftarrow(L',R')$
and $(\ell,r)\leftarrow(L',R')$.
\item If $a$ is \emph{halt}, the run halts with no output.
\item If the new state has occurred earlier in the run, the run halts with no output.
\end{enumerate}
\end{definition}

\begin{definition}[Output function]\label{def:output}
$\delta_{\rho,n}:\Gamma^{n}\to D_{n}\cup\{\bot\}$, where $D_{n}=\{i/(n-1):0\le i\le n-1\}$ and
$\bot$ denotes ``no output'', is the map defined by Definition~\ref{def:run}. It is total: the state
space is finite, so every run halts by clause 5.
\end{definition}

\begin{remark}[The repeat test costs no space]\label{rem:repeat}
Clause 5 is a property of the orbit, not a resource of the device: $f$ and the geometry determine the
next state, so a device can implement the test by Brent's algorithm in $O(1)$ extra work bits. It
costs reads and steps, never space. The clause therefore does not smuggle an $\Omega(n^{2})$-bit
table into the model --- but a device that merely simulates must pay for cycle detection, and the
read counts of this paper include that payment.
\end{remark}

\subsubsection{Two readings of $\varphi$, and which is in force where}

\begin{definition}[Oracle semantics, (O)]\label{def:oracle}
$\varphi$ is an oracle: evaluating $\varphi(w)$ costs $0$ reads and $0$ steps, for every window $w$.
\end{definition}

\begin{proposition}[The oracle reading is vacuous]\label{prop:oracle}
Under (O), every rule $\rho$ costs zero reads, so every rule is trivially fully linear and none is
superlinear.
\end{proposition}

\begin{proof}
The device stores $(L,R,\ell,r)$ --- four integers below $n$, so $4\lceil\log_{2}n\rceil$ bits ---
queries the oracle for $\varphi(\tau[\ell..r])$, looks $f$ up in its finite control, and updates. It
never senses a cell, so reads $=0$. Clauses 1--5 of Definition~\ref{def:run} are decided from the
stored state and the oracle's answer.
\end{proof}

\begin{definition}[Computed semantics, (C)]\label{def:computed}
The device has no oracle. On input $\tau$ it must write $\delta_{\rho,n}(\tau)$ on its output tape,
sensing only cells of $\tau$. $\varphi$ enters only through $\delta_{\rho,n}$; the device may compute
$\varphi$ however it likes, precompute it, cache it, or bypass it entirely.
\end{definition}

\begin{definition}[$\varphi$-state; incremental computability]\label{def:phistate}
A \emph{$\varphi$-state} for $\rho$ on a window $w$ is a value $\sigma\in\Sigma$ together with an
evaluation map $\mathrm{ev}:\Sigma\times\mathbb N\to A$ such that $\mathrm{ev}(\sigma,|w|)=\varphi(w)$.
A device \emph{carries} $\varphi$ if its work tape holds a $\varphi$-state and $\mathrm{ev}$ costs
$O(1)$ steps and $0$ reads. $\varphi$ is \emph{incrementally computable} if there is
$\mathrm{upd}:\Sigma\times\Gamma\to\Sigma$, costing $O(1)$ steps and $0$ reads, with
$\mathrm{upd}(\sigma_{w},c)=\sigma_{w\setminus c}$ for every window $w$ and every cell $c$ at either
end of $w$.
\end{definition}

\begin{remark}[Which reading is in force where]\label{rem:which}
The three readings occur in different parts of the paper, and the classification is only meaningful
under (C).
\begin{center}
\small
\begin{tabular}{@{}p{.26\textwidth}p{.14\textwidth}p{.50\textwidth}@{}}
\toprule
location & reading & what it establishes \\
\midrule
Definition~\ref{def:rule} & (O) & the \emph{semantics} of the rule: $f$ is handed
$\varphi(\tau[\ell..r])$. Not a cost statement.\\
the machines of Section~\ref{sec:cost} & (C) & read counts of devices that compute $\varphi$
themselves; every number in the tables is of this kind.\\
Theorem~\ref{prop:collapse}, Theorem~\ref{thm:rank} & (C) with a carried $\varphi$-state & the only
places where the assumption is written down; strictly stronger than Definition~\ref{def:rule}.\\
the ceiling of Section~\ref{sec:ceiling} & neither & a statistical estimator with a read budget, not
required to be exact; not a device in the sense of Definition~\ref{def:rule}.\\
\bottomrule
\end{tabular}
\end{center}
\end{remark}

\noindent The consequence of Proposition~\ref{prop:oracle} is worth stating plainly: a paper that
counted reads under (O) would classify nothing, since every rule would cost zero. Every read count
below is therefore a statement under (C), and the two theorems that mention a carried $\varphi$-state
are the only places where that assumption is used at full strength.

\subsection{The predictor, and $\kappa$}\label{app:kappa}

The classification uses $\kappa$, and $\kappa$ is a supremum, so it needs both a definition and a
statement of what a measurement of it can and cannot establish.

\begin{definition}[Predictor]\label{def:predictor}
Fix $\rho$, $n$ and a budget $B=B(n)$. A \emph{predictor} is a device that, on input $\tau$, reads at
most $B(n)$ cells of $\tau$ --- adaptively, in any order --- performs arbitrary computation on what it
has read, and outputs one value in $D_n$. It is not required to be correct; it is scored. The tapes
that concern it are the \emph{meeting} tapes, those on which $\delta_{\rho,n}$ is defined: a tape with
no output has no $\delta$ to predict and is not counted.
\end{definition}

\begin{definition}[Score, and $\kappa$]\label{def:kappadef}
For a predictor $P$ let $R^2_n(P)$ be the coefficient of determination of $P$'s output against
$\delta_{\rho,n}$, computed over the meeting tapes with cross-validation and with the constant
predictor in the comparison. Define
\[
\kappa_{\rho}(c)=\lim_{n\to\infty}\ \sup_{P\,:\,B=cn} R^2_n(P),
\qquad
\kappa_{\rho}=\sup_{c\ge1}\kappa_{\rho}(c).
\]
The classes are read in this order: \emph{fully linear} if $g_{\rho}=O(n)$; otherwise
\emph{semi-linear} if $e_{\rho}=O(n)$; otherwise \emph{superlinear}.  The superlinear class is
not split further: $\kappa_{\rho}$ measures how close a rule comes to a cheap exact machine, no
member of the catalogue reaches $\kappa_{\rho}=1$, and Conjecture~\ref{conj:perfect} says none can.
\end{definition}

\begin{remark}[What a measurement of $\kappa$ can establish]\label{rem:kappasup}
$\kappa$ is a supremum over predictors and a limit in $n$, so a measurement is a \emph{lower} bound:
it exhibits one predictor and reports what it achieved. Two consequences follow.
\begin{enumerate}
\item $\kappa=1$ \emph{can} be established, by exhibiting a device that reads at most $cn$ cells and
reproduces $\delta$ on every meeting tape.
\item $\kappa=0$ \emph{cannot} be established by measurement. The strongest statement a table can
carry is that every predictor tried failed to beat the constant predictor.
\end{enumerate}
\end{remark}

\begin{proposition}[Linear mean forces perfect prediction]\label{prop:elinear}
If $e_{\rho}=O(n)$ then $\kappa_{\rho}=1$.
\end{proposition}

\begin{proof}
Let $M$ be a device with mean cost $e\le Cn$ over the meeting tapes. Cap it at $cn$ reads and let it
guess on the tapes where the cap stops it. By Markov's inequality the capped fraction is at most
$C/c$, and since $\delta\in[0,1]$ the squared residual on those tapes is at most one, so
\[
R^2_n(c)\ \ge\ 1-\frac{C}{c\,\mathrm{Var}(\delta)}
\]
for every $n$ and every $c$. Taking $n\to\infty$ and then $c\to\infty$ gives
$\kappa_{\rho}=1$.
\end{proof}

\noindent The contrapositive is what the classification uses: $\kappa_{\rho}<1$ implies
$e_{\rho}=\omega(n)$, so every rule with $\kappa_{\rho}<1$ is superlinear. The converse
fails --- a rule with superlinear mean cost can still be asymptotically predictable --- and that
failure is exactly the boundary at $\kappa=1$.

\begin{remark}[$R^2$ is blind to one-cell errors]\label{rem:grid}
$\delta$ takes values on a grid of spacing $1/(n-1)$. A predictor that is off by exactly one grid cell
on a constant fraction of tapes still has $R^2\to1$, because the rescaled residual is
$(n-1)\,|\hat\delta-\delta|=1$ and its square is $1/n^2$ against a total variance of order one. A
rising $R^2$ therefore does not show that the misses are shrinking. The quantity to report is the
residual \emph{in grid units}: the fractions with $(n-1)|\hat\delta-\delta|$ equal to $0$, to $1$ and
to at least $2$, and how those fractions move with $n$.
\end{remark}

\subsection{The catalogue of rules}\label{app:catalogue}

Every rule the main text names appears in Table~\ref{tab:rule-catalogue}, one row each, with its
$\varphi$ written on a single line precise enough to reimplement. Write $m=r-\ell+1$,
$k=\max(1,\lfloor m/2\rfloor)$ and $a_i=2\tau_i\in\{0,1,2\}$; let $s$ be the label tape and $G$ its
prefix walk, $G_0=0$, $G_{j+1}=G_j+s_j$. Except where a row states otherwise,
$\varphi=(b_2,b_1)$ with the concavity test
$b_2=b_2^{\mathrm{conc}}=[\,k\sum_{i=0}^{m-1}a_{\ell+i}>m\sum_{i=0}^{k-1}a_{\ell+i}\,]$ as its second
coordinate, so each cell of the table records $b_1$ alone. The rows marked four-valued read the label
tape $s$; the three-valued rows read only $\sigma=a$; by the embedding lemma every $\sigma$-rule is
bit-identical in the four-valued model, so the two columns are a statement about what the rule needs,
not about which engine ran it.

The action convention is the canonical four-move table --- $\ell\mapsto\ell-1$ on $(0,0)$,
$r\mapsto r+1$ on $(0,1)$, $r\mapsto r-1$ on $(1,0)$, $\ell\mapsto\ell+1$ on $(1,1)$, the headers
moving inward exactly when the two bits agree --- unless the row says otherwise, and three rows do.
\texttt{hashmod} uses a two-action map that ignores $b_2$: $b_1=0$ moves $r\mapsto r-1$ and $b_1=1$
takes a left trim. \texttt{PC} and \texttt{CX} use a three-clause convention on $(m,W)$ with
$W=R-L+1$: move $\ell\mapsto\ell+1$ when $m\le2$, move inward on $b_1$ when $W-m<k(W)$ with
$k(W)=\max(1,\lfloor W/64\rfloor)$, and otherwise trims on the side fixed by the parity of $\ell-L$.
Two exit-trim conventions occur in the record, and they differ whenever the window does not touch the
boundary it leaves through. The canonical engines trim the side \emph{opposite} to the one the move
would have left; the sign-blind and the four-valued engines trim the \emph{same} side. Only the
canonical row carries a dagger below and uses the first; every other four-move row uses the second.

\begin{table}[p]
\centering
\footnotesize
\setlength{\tabcolsep}{4pt}
\renewcommand{\arraystretch}{1.02}
\makeatletter
\newcommand{\catrr}{\raggedright\let\\\@arraycr}
\makeatother
\begin{tabular}{@{}p{0.185\textwidth}lp{0.59\textwidth}l@{}}
\toprule
rule & tape & $\varphi$ & class \\
\midrule
\catrr \texttt{canonical}$^{\dagger}$ & 3-valued &
\catrr $\big[\,\widehat{\mathrm{corr}}_1(a_\ell,\dots,a_r)>0\,\big]$, with $[a_\ell>\tfrac12]$
at $m{=}1$, $[a_r>a_\ell]$ at $m{=}2$, and $0$ when either lagged half is constant &
FL \\
\catrr \texttt{monotone family} & any &
\catrr none: the family is defined by $f$, not by $\varphi$ & FL \\
\catrr \texttt{NBLK} (statistic) & 4-valued &
\catrr $\mathrm{NBLK}(w)=\#\{\,i\in(\ell,r+1]\;:\;G_i=G_\ell\,\}$ --- the returns of the
label walk to its starting level & --- \\
\catrr \texttt{NBLK=1} & 4-valued &
\catrr $\big[\mathrm{NBLK}(w)=1\big]$ & SU \\
\catrr \texttt{NBLK>2} & 4-valued &
\catrr $\big[\mathrm{NBLK}(w)>2\big]$, count $\ge3$ & SU \\
\catrr \texttt{NBLK>3} & 4-valued &
\catrr $\big[\mathrm{NBLK}(w)>3\big]$, count $\ge4$ & SU \\
\catrr \texttt{AMAX} & 4-valued &
\catrr $\big[G_{r+1}>\max_{\ell\le i\le r}G_i\big]$, the window's border attains a strict
maximum & SU \\
\catrr \texttt{AMIN} & 4-valued &
\catrr $\big[G_{r+1}<\min_{\ell\le i\le r}G_i\big]$ & SU \\
\catrr \texttt{MAXP} & 4-valued &
\catrr $\big[(\max_{\ell\le i\le r+1}G_i-G_\ell)^2>m\big]$ & SU \\
\catrr \texttt{KCHG} & 4-valued &
\catrr $c=\#\{p\in[\ell+1,r]:s_{p-1}s_p<0\}$, $p_c=$ the $\lceil c/2\rceil$-th such $p$;
$\big[2(p_c-\ell)>m\big]$, and $0$ if $c=0$ & SU \\
\catrr \texttt{LAS} & 4-valued &
\catrr $\mathrm{LAS}(w)=1+\#\{\text{sign changes in }(s_i)_{i\in w,\,s_i\ne0}\}$, (0 if none); $\big[4\,\mathrm{LAS}>m\big]$ & SU \\
\catrr \texttt{BAL>m/2} & 4-valued &
\catrr $\mathrm{BAL}(w)=\max\{j-\ell:\ell<j\le r+1,\ G_j=G_\ell\}$, (0 if none);
$\big[2\,\mathrm{BAL}>m\big]$ & SU \\
\catrr \texttt{BAL>m/4} & 4-valued &
\catrr same $\mathrm{BAL}$; $\big[4\,\mathrm{BAL}>m\big]$ & SU \\
\catrr \texttt{hashmod} & 3-valued &
\catrr $\big[H(w)\bmod m=0\big]$, $H(w)=\big(\sum_{j=0}^{m-1}2^j a_{\ell+j}\big)\bmod p$,
$p=10^{9}+7$ & SU \\
\catrr \texttt{dense2} & 3-valued &
\catrr $\big[4\,\#\{i\in[\ell,r]:a_i=2\}>m\big]$ & SU \\
\catrr \texttt{runs} & 3-valued &
\catrr $\big[2\,\#\{i\in[\ell,r-1]:a_i\ne a_{i+1}\}>m\big]$ & SU \\
\catrr \texttt{pal} & 3-valued &
\catrr $\big[a_\ell a_{\ell+1}\cdots a_r=a_r\cdots a_{\ell+1}a_\ell\big]$ & FL \\
\catrr \texttt{pattern} & 3-valued &
\catrr $\big[\exists j\in[\ell,r-2]:(a_j,a_{j+1},a_{j+2})=(2,1,1)\big]$, and $0$ if $m<3$ &
FL \\
\catrr \texttt{longrun} & 3-valued &
\catrr $\big[4\,\Lambda_2(w)>m\big]$, $\Lambda_2(w)=\max\{t:\exists j,\ a_j=\cdots=a_{j+t-1}=2\}$ & SU \\
\catrr \texttt{peak} & 3-valued &
\catrr $\big[6\,\#\{i\in[\ell+1,r-1]:a_{i-1}<a_i>a_{i+1}\}>m\big]$ & SU \\
\catrr \texttt{trip} & 3-valued &
\catrr $\big[6\,\#\{i\in[\ell,r-2]:a_i=a_{i+1}=a_{i+2}\}>m\big]$ & SU \\
\catrr \texttt{oddrun} & 3-valued &
\catrr $\big[2\,\#\{\text{maximal constant runs of }w\text{ of odd length}\}>m\big]$ &
SU \\
\catrr \texttt{ballast} & 3-valued &
\catrr the monotone collapse of Section~\ref{sec:classes} with the one-shot defect defined
there & SL \\
\catrr \texttt{PC} & 4-valued &
\catrr $p_0=0$, $p_{t+1}=(3p_t+1+a_{\ell+p_t})\bmod m$, $\big[2p_m\ge m\big]$ & --- \\
\catrr \texttt{CX} & 4-valued &
\catrr $q=2^{31}-1$, $h(k)=3^{k^2+1}\bmod q$,
$B(w)=\sum_{i=\ell}^{r}a_i\,h(i-\ell)\bmod q$, $\big[2B<q\big]$ & --- \\
\bottomrule
\end{tabular}
\caption{The catalogue. Each row is one rule with the tape it reads, its $\varphi$ in one line, and
the class it is placed in: FL fully linear ($g_\rho=O(n)$), SL semi-linear ($e_\rho=O(n)$ with
$g_\rho=\omega(n)$), SU superlinear ($e_\rho=\omega(n)$), --- unclassified. The superlinear rows are
not subdivided here: their measured $\kappa_\rho$ is in Table~\ref{tab:kappa-inc}, and every one of
them is strictly between $0$ and $1$. A dagger marks the one row that uses the opposite-side exit
trim; row~1 of the table is the only such row.}
\label{tab:rule-catalogue}
\end{table}

Three of these rows --- \texttt{BAL}, \texttt{KCHG} and \texttt{LAS} --- have no prose definition
anywhere in the record: the manuscript, the measurement notes and the appendix all name them and
none defines them. Their $\varphi$ above is taken from the implementation rather than from a stated
definition, so a reader who wants to check them should read the code, and should expect the
definition it finds to be the one that produced the numbers rather than the one a gloss would
suggest. \texttt{LAS} is the sharpest case: as implemented it coincides with \texttt{runs}, which is
why its separate classification rests on measurement alone.

\subsection{Structural facts, and what of them is proved}\label{app:proofs}

Each item below is labelled \emph{proved} or \emph{verified}. The distinction is not decorative: three
of the six are verified on exhaustive small cases and are not theorems, and saying so is cheaper than
having a referee say it.

\paragraph{Parity of the boundary (proved).} Write $C$ for the number of cells of the window exceeding
its mean and $D=2C-m$ for the signed surplus. Since $2C\equiv m \pmod 2$, $D\equiv m\pmod 2$, so the
boundary $D=0$ --- the only state in which a scan can change an epoch's outcome --- is reachable only
at even $m$. Verified on all $27\,984$ epochs of the cube for $n=4$--$8$: $b_2=0$ forces a double trim
without exception, and every epoch whose outcome differs from the $b_1$-side trim visits $D=0$.

\paragraph{Drift (proved for even $n$).} For the ensemble of Definition~\ref{def:rule},
$\mathbb E[4N_2-m]=0$ at even $n$ and $+2$ at odd $n$, where $N_2$ counts the cells equal to $2$. Hence
\texttt{dense2}'s bit has no drift at even $n$, flips in $35$ per cent of epochs, and flips the trim
side with it. Measured consequence: enlarging the endgame buffer from $24$ to $320$ cells leaves the
error rate at $14$ per cent.

\paragraph{Level counts for the extremum rules (verified).} For a $\pm1$ walk, the statement ``the
endpoint $r$ is the strict maximum of the window'' is equivalent to ``the level $G_{r+1}$ is never
visited in $[\ell,r]$''. The equivalence turns \texttt{AMAX} and \texttt{AMIN} into single-level
counts, which is what makes them reachable by the same band$+$frontier machine as \texttt{NBLK}.
Verified on $132\,036$ windows with no mismatch. This is a verified equivalence, not a proved one; the
converse direction is the one that could fail on a longer window, and it was checked only on the
windows the walk visits.

\paragraph{The rank of the convolution kernel (proved).} For $h(k)=3^{k^{2}}$ and modulus $q$,
\[
H_{m}[d,k]=h(k+d)=3^{(k+d)^{2}}=3^{d^{2}}\cdot 9^{dk}\cdot 3^{k^{2}},
\]
so $H_{m}=3\,\mathrm{diag}(3^{d^{2}})\,[9^{dk}]\,\mathrm{diag}(3^{k^{2}})$. The middle factor is a
Vandermonde matrix on the nodes $9^{d}$, $d=0,\dots,m-1$, which are pairwise distinct whenever
$\mathrm{ord}_{q}(9)>m-1$. With $q=2^{31}-1$, the order is $357\,913\,941$ --- computed, not assumed:
$q$ is prime, the order divides $q-1$, and it exceeds the largest window the tables use by a factor of
$21\,845$ (\texttt{probe\_vandermonde/vand\_check.py}) --- so the nodes are pairwise distinct for every
$m$ in the range studied and $\operatorname{rank}H_{m}=m$ exactly over $\mathbb F_{q}$. The two diagonal
factors are powers of $3$, nonzero in $\mathbb F_{q}$ and hence invertible, so the rank is that of the
Vandermonde factor. This is the only rank computation the paper
relies on.

\paragraph{The witness family for the cubic cost (verified).} The family is $1^{m}$ --- the constant
tape --- with nine cells changed. The machine's two encodings cover up to eight defects and it uses nine,
which its accumulator still absorbs; the cost per test is flat in the window length for every fixed
count that was checked, up to sixteen. On this family the Brent-style machine costs $n^{3}/685$ at $n=4096$ with local
exponent $2.90$, and $C(m)/m$ is flat in $m$ (per-window exponent $1.07$ on the artificial family,
$0.95$ on the real tape). The dependence on the defect \emph{count} rather than density is what makes
the exponent $3$ rather than $2.5$; with a fixed defect density the same machine gives $2.5$ to
$2.85$. This is a measured witness, not a lower bound: no argument in this paper shows that no better
machine exists.

\paragraph{Why one number suffices (proved).} If $b_1$ is constant over the epoch's range and that
constancy is certified by a tape-global maximum, the epoch's effect is a single trim whose side is a
function of $(b_2,b_1)$ evaluated once, so the strategy costs $O(n\rho)$, where $\rho$ is the longest
run of $2$s --- this is \texttt{longrun}. If the certificate is instead a last-passage query on a
$\pm1$ walk, its exact structure needs $\Theta(n)$ words and the strategy does not linearise. The
dichotomy is the reason the certificate's \emph{shape}, not its size, decides the class.

\paragraph{Step counts, and why there is no clock}\label{app:steps}
The body states no rule-independent clock, because the measurement contradicts it. On one ensemble,
with every run allowed to finish, the mean step count at $n=4096$ divided by $n^{2}/128$ is:

The ratio ranges from $0.385$ for \texttt{NBLK>3} to $15.58$ for \texttt{hashmod}, a factor of
forty, and the growth exponent between $n=1024$ and $4096$ ranges from $1.70$ to $2.20$. These
are the step counts of the machines of Section~\ref{sec:cost}, regenerated from their traces;
an earlier draft printed a wider span whose low end was not in the data. A clock would have to
be a property of the model; this is a property of the rule.

\paragraph{The collapse (verified, not proved).} The collapse (Proposition~\ref{prop:collapse}) was
checked on $320\,880$ queries with no error, and the body states it as a proposition for that
reason; the general-$n$ statement is Conjecture~\ref{conj:collapse}.

\paragraph{The general bound, and why it is not sharp.} In cell reads the trivial bound is $O(n^{4})$:
the walk makes $O(n^{2})$ steps and a step that re-evaluates a window statistic costs up to $O(n)$
reads. Under the epoch structure the bound is $O(n^{3})$, and the five level-count rules of
Section~\ref{sec:cost} meet the next rung down at $n^{2}$ with the band$+$frontier machine. The
$n^{3}$ rung is occupied by \texttt{PC} on the witness family above. What is missing is not a better
machine but a lower bound: Section~\ref{sec:ceiling} records the two ceilings that block
the standard arguments.
\subsection{Protocols and gates}
\label{app:protocols}

Every machine is written read-exact and gated before any timing, at zero tolerance: the walk against
the literal definition on the full cube $\{0,1,2\}^{n}$ for $n\le10$ ($88\,560$ tapes) and on random
tapes at $n=11$--$96$, accepted only at zero mismatches in $\delta$ \emph{and} in the step count; the
$b_1$ table against an independent implementation on $2\,192\,103$ windows; and the registers after
every trim against a prefix-sum oracle. Three bugs in our own machines were caught this way after their
numbers had been computed once, and the numbers were recomputed.

Two protocol choices are recorded because each changed a number. Exponents are fitted on the largest
three sizes only: fitting from $n=64$ produced two exponents we later withdrew ($1.696$ for
\texttt{peak} and $1.826$ for \texttt{trip}; both are $2.04$--$2.12$ on the large sizes), because the
$2n$ initialisation dominates at small $n$. And the worst case is measured on a constructed family
rather than by hill climbing: on \texttt{NBLK>2} a climb found $4.3\,n^{2}$ where the family found
$48.9\,n^{2}$ at the same size, a factor of eleven.

\subsection{Verification gates}
\label{app:gates}

Every machine used for a number in this paper passed a gate of the
following form before any timing. The counts are totals over the paper's
strategies.

\begin{table}[h]
\centering
\small
\begin{tabular}{@{}lll@{}}
\toprule
what is checked & how & count \\
\midrule
the walk against the literal definition & full cube $n\le10$ & $88\,560$ tapes \\
the same, four-valued model & cube $+$ support $T_n$ & $130\,500$ walks \\
the machine's registers after every trim & prefix-sum oracle & every scan step \\
$b_1$ on all windows of the cube & independent implementation & $2\,192\,103$ windows \\
the endgame at $W=n$ reduces to the walk & cube $n\le10$ & $147\,555$ walks \\
the constructed worst-case family is realisable & by construction & all members \\
\bottomrule
\end{tabular}
\caption{Gates. A number is reported only if its machine passes. Three real
bugs in our own machines were caught this way after the numbers had
already been computed once, and the numbers were recomputed.}
\end{table}

\subsection{Status of each claim}
\label{app:status}

\begin{table}[h]
\centering
\small
\begin{tabular}{@{}p{0.46\textwidth}p{0.20\textwidth}p{0.24\textwidth}@{}}
\toprule
claim & status & where \\
\midrule
parity of the $D=0$ boundary & proved & Appendix~\ref{app:proofs} \\
\texttt{dense2} has no drift at even $n$ & proved & Appendix~\ref{app:proofs} \\
$g_\rho=O(n^4)$ in cell reads, $O(n^3)$ under the epoch structure & proved & Appendix~\ref{app:proofs} \\
$g_\rho=\Theta(n^3)$ for machines that recount per test & proved & Appendix~\ref{app:proofs} \\
global-max certificate $\Rightarrow$ $O(n\rho)$ & proved & Appendix~\ref{app:proofs} \\
\texttt{longrun}: $e/n$ rises $8.53\to9.39$ over $n=512\to16384$, $\alpha=1.024$ & measured & Section~\ref{sec:classes} \\
\texttt{runs} costs $2.77\,n$ on the CES device, exponent $1.00$ & measured & Section~\ref{sec:cost} \\
$\kappa\in[0.472,0.925]$ for thirteen rules at $c=16$ & measured & Section~\ref{sec:kappa} \\
$\kappa_{\texttt{hashmod}}\approx0$ & measured & Section~\ref{sec:kappa} \\
\midrule
no strategy has $g_\rho=\omega(n^3)$ & \textbf{open} & --- \\
every rule outside the fully linear class has $g_\rho=\omega(n)$ & \textbf{open} & --- \\
\midrule
``every non-linear rule has $g_\rho=\Omega(n^2)$'' & \textbf{false} & the constructed
$n^{3/2}$ rung ($g=\Theta(n^{1.5})$) and \texttt{ballast} ($g=\Theta(n^2)$) \\
some rule costs $n^{3}/685$ on the witness family & measured, upper bound & Section~\ref{sec:cost} \\
``a window statistic can be evaluated in $O(1)$ reads'' & \textbf{false in general} & Appendix~\ref{app:proofs} \\
\bottomrule
\end{tabular}
\caption{What is proved, what is measured, and what is open. ``Measured''
means the protocol of Appendix~\ref{app:protocols} ran and the gate of
Appendix~\ref{app:gates} passed; it does not mean a lower bound was
proved.}
\end{table}

\subsection{The ensemble}
\label{app:ensemble}

A tape of length $n$ is drawn by choosing, independently for each of two
columns, a uniform $h$-subset of $\{1,\dots,n\}$ with $h=\lceil n/2\rceil$;
a cell's value is the sum of its two bits, so values lie in $\{0,1,2\}$
with probabilities $(1/4,1/2,1/4)$ up to the dependency induced by the
subsets. All averages are over this ensemble. Strategies read only
$\sigma=b_1+b_2$ except where the text says the label $s=b_1-b_2$ is used.
